\documentclass{article}
\usepackage[margin=1.15in]{geometry}
\usepackage{xcolor}
\usepackage{parskip}
\usepackage{fancyhdr}
\usepackage{array}
\usepackage{booktabs}
\usepackage{enumitem}
\usepackage{amsmath}
\usepackage{tabularx}

\pagestyle{fancy}
\fancyhf{}
\cfoot{\thepage}
\rhead{Lecture 4}
\lhead{MA 214: Applied Statistics (Spring 2026)}
\renewcommand{\headrulewidth}{.4mm}

\newcommand{\blankline}[1]{\underline{\hspace{#1}}}

\begin{document}

\noindent\textbf{Name:}\blankline{2.25in}\hfill\textbf{BUID:}\blankline{1.55in}\newline

\begin{center}
 \Large In-Class Worksheet: Using Regression Models
\end{center}

\subsection*{Scenario}
Today we practice using a fitted regression model after it has been fit: interpreting correlation and $R^2$, making predictions, checking for extrapolation, and using a 0/1 indicator variable for a categorical predictor.

\medskip
\noindent\textbf{Goals:} \textbf{(1)} interpret $r$ and $R^2$, \textbf{(2)} compute predicted values and residuals, \textbf{(3)} identify interpolation vs. extrapolation, \textbf{(4)} interpret an indicator-variable model.

\subsection*{Part 1: Correlation and $R^2$}
Suppose a simple linear regression model uses high school graduation rate to predict poverty rate across states. The correlation between the two variables is $r=-0.75$.

\begin{enumerate}[label={\bfseries (\Alph*)},itemsep=.8em]
\item Describe the direction and strength of the linear relationship. \\[2.5em]
\item Compute $R^2$.
\[
R^2 = \blankline{2.2in}
\]
\item Interpret $R^2$ in context. \\[3.5em]
\item Which statement is the best interpretation of $R^2=0.56$? Circle one.
\begin{enumerate}[label=\arabic*.,itemsep=.2em]
\item The model predicts 56\% of states exactly correctly.
\item HS graduation rate causes 56\% of the poverty rate.
\item 56\% of the variability in poverty rate is explained by the model.
\item The correlation is 0.56.
\end{enumerate}
\end{enumerate}

\subsection*{Part 2: Prediction Checklist}
Use the fitted model
\[
\widehat{\text{poverty}} = 64.78 - 0.62(\text{HS grad}),
\]
where poverty rate and HS graduation rate are both measured in percent.

\begin{enumerate}[label={\bfseries (\Alph*)},itemsep=.8em]
\item Predict the poverty rate when HS graduation rate is 90\%.
\[
\widehat{\text{poverty}} = \blankline{4in}
\]
\item Suppose the observed poverty rate is 7.5\%. Compute the residual.
\[
e = y - \hat y = \blankline{3.5in}
\]
\item Is the observed point above or below the fitted line? Explain using the sign of the residual. \\[3em]
\item Should the actual poverty rate equal the predicted value exactly? Why or why not? \\[3em]
\end{enumerate}

\subsection*{Part 3: Interpolation or Extrapolation?}
Suppose the observed HS graduation rates in the data range from 77\% to 92\%.

\begin{center}
\renewcommand{\arraystretch}{1.45}
\begin{tabularx}{\textwidth}{p{4.2cm}p{4.1cm}X}
\toprule
\textbf{Prediction at} & \textbf{Interpolation or extrapolation?} & \textbf{Main concern} \\
\midrule
85\% HS graduation rate & & \\[1.7em]
93\% HS graduation rate & & \\[1.7em]
65\% HS graduation rate & & \\
\bottomrule
\end{tabularx}
\end{center}

\begin{enumerate}[label={\bfseries (\Alph*)},itemsep=.8em]
\item A fitted model uses house sizes from 800 to 3500 square feet. Which prediction is most risky? Circle one and explain briefly.
\begin{center}
1200 square feet \qquad 2600 square feet \qquad 9000 square feet
\end{center}
\vspace{3em}
\end{enumerate}

\subsection*{Part 4: A Categorical Predictor}
Suppose we fit a model for poverty rate using region, where
\[
\text{south} =
\begin{cases}
1, & \text{Southern state},\\
0, & \text{otherwise}.
\end{cases}
\]
The fitted model is
\[
\widehat{\text{poverty}} = 10.19 + 3.47(\text{south}).
\]

\begin{enumerate}[label={\bfseries (\Alph*)},itemsep=.8em]
\item What is the reference group? \\[2.5em]
\item What is the predicted poverty rate for a non-Southern state?
\[
\blankline{2.5in}
\]
\item What is the predicted poverty rate for a Southern state?
\[
\blankline{2.5in}
\]
\item Interpret the coefficient 3.47 in context. \\[3.5em]
\end{enumerate}

\subsection*{Wrap-Up}
Write one sentence explaining why a prediction can be mathematically easy but statistically risky. \\[4em]

\end{document}
